\SourceLeaf{32}oder richtiger, dass jene Summe niemals positiv werden könne. Sind daher $\gamma,\gamma',\gamma''$ u. s. w. von $c,c',c''$ u. s. w. verschiedene, aber mit den Bedingungen des Systems verträgliche Plätze; und $\theta,\theta';\theta''$ u. s. w. die Winkel, welche $c\gamma,c'\gamma',c''\gamma''$ u. s. w. mit $cb,c'b',c''b''$ u. s. w. machen, so ist allemal $\Sigma m.cb.c\gamma.\cos\theta$ entweder $0$ oder negativ. Da nun
\[
\gamma b^2=cb^2+c\gamma^2-2cb.c\gamma.\cos\theta
\]
so ist klar, dass
\[
\Sigma m.\gamma b^2-\Sigma m.cb^2=\Sigma m.c\gamma^2-2\Sigma m.cb.c\gamma.\cos\theta
\]
folglich immer positiv sein wird, also $\Sigma m.\gamma b^2$ immer grösser als $\Sigma m.cb^2$, d. i. dass $\Sigma m.cb^2$ ein Minimum sein wird. W. Z. B. W.

Es ist sehr merkwürdig, dass die freien Bewegungen, wenn sie mit den nothwendigen Bedingungen nicht bestehen können, von der Natur gerade auf dieselbe Art modificirt werden, wie der rechnende Mathematiker, nach der Methode der kleinsten Quadrate, Erfahrungen ausgleicht, die sich auf unter einander durch nothwendige Abhängigkeit verknüpfte Grössen beziehen. Diese Analogie liesse sich noch weiter verfolgen, was jedoch gegenwärtig nicht zu meiner Absicht gehört.
\par\medskip\centerline{\rule{.3\textwidth}{.3pt}}
